\chapter{Introduction}

 The introduction sets the stage for the dissertation by presenting the research context, defining the problem, engaging readers, and providing a foundation for subsequent chapters. Compared to the abstract, it is more comprehensive, offering a detailed overview of the study’s background and significance.

\section{Content Elements}
\begin{itemize}
    \item Nature of the Research Problem: Clearly articulate the problem, its context, and its academic or practical importance.
    \item Literature Review: Summarize and critique prior studies, highlighting progress, limitations, and research gaps.
    \item Methodology Overview: Briefly describe the methods, procedures, or tools used, emphasizing their novelty or improvements over existing approaches.
    \item Key Results and Conclusions: Summarize major findings and their implications.
    \item Research Impact: Explain the study’s contributions to academia or practice.
\end{itemize}

\section{Writing Tips}
\begin{itemize}
    \item Begin with an engaging hook, such as a question, phenomenon, or compelling data point.
    \item Focus the literature review on directly relevant studies, organizing by theme or chronology to avoid a laundry-list approach.
    \item Use a problem statement or list of research questions to sharpen the focus for readers.
    \item Conclude with an overview of the dissertation’s structure to guide readers through the document.

\end{itemize}
